\section{The height-four obstruction calculation}
\label{may:section}

The algebraic spectral sequences below use the filtered cobar complexes
of Section~\ref{may:preliminaries}.

\subsection{Possible differential targets}

We now use the finite spectrum $Y$ furnished by
Proposition~\ref{low:V3}.  Its homology is $R/I_4$, its bottom cell is a
unit, and it has a binary two-sided unital multiplication.  The possible
positive Adams--Novikov differentials from $v_4$ have targets
\begin{equation}
 s=qk+1,\qquad w=u_4+k,\qquad k\geq1.
 \label{may:obstruction-degrees}
\end{equation}

\begin{lemma}[Explicit cutoffs]
\label{may:target-bound}
The groups in \eqref{may:obstruction-degrees} are identified with
$H^{s,w}(P)$.  Their May page has only positions $0,1,2,3$, and it is zero
unless $k\leq p+2$.  For $k\geq2$, only positions $0,1,2$ can occur.
\end{lemma}

\begin{proof}
If $w<s$, the degree-$(s,w)$ cochain groups for both $Y$ and $P$
are zero, as is the May page, and all the conclusions are immediate.
We may therefore suppose $w\geq s$.  The normalized bound in
Lemma~\ref{fw:normalized-bound} then gives
$(q-1)k\leq u_4-1$.  Consequently
\[
 w\leq u_4+\frac{u_4-1}{q-1}
 <\frac q{q-1}u_4<4p^3<p^4
 \quad (p=1009),
\]
where $u_4<2p^3$ and $q/(q-1)<2$.  Thus no letter can reach position
four.  Also $w<u_5$.  A coefficient containing $v_4$ leaves at most $k$
weight for normalized bars, whereas the three lengths $s-1,s,s+1$ are
all greater than $k$, because $s-1=qk>k$.  Higher coefficient generators
are excluded by $w<u_5$.  The adjacent-chain comparison of
Lemma~\ref{fw:additive-comparison} therefore identifies the required
cohomology with that of $P$.

There are ten possible exterior intervals on positions $0,1,2,3$.
Since $s$ is odd, $m$ is odd and hence $m\leq9$.  Every polynomial letter
has weight at least $p$, so
\[
 u_4+k\geq pL\geq p((p-1)k-4),
 \qquad
 (p^2-p-1)k\leq u_4+4p.
\]
But
\[
 (p^2-p-1)(p+3)-(u_4+4p)=p^2-9p-4>0
 \quad (p=1009).
\]
It follows that $k\leq p+2$.

For $k\geq2$ we have $L\geq2p-6$.  If any letter reaches position
three, that letter has weight at least $p^3$, while all other polynomial
letters have total weight at least $(L-1)p$.  This is a valid lower bound
whether the distinguished letter is exterior or polynomial.  Thus
\[
 w\geq p^3+2p^2-7p
 >p^3+p^2+2p+3=u_4+p+2\geq w,
\]
since the strict difference is $p^2-9p-3>0$ at $p=1009$.  This
contradiction excludes position three in that range.
\end{proof}

\subsection{The first obstruction}

The first possible differential is $d_{q+1}$, and its target is the
actual cobar cohomology in degree $(q+1,u_4+1)$. We first explain why
an injective multiplication on this group will rule out the differential.
Lemma~\ref{fw:v5-relation} gives
\[
 v_4h_{1,4}+e v_4^p h_{1,0}=0
\]
in the $E_2$ algebra for $Y$. By $q$-sparsity no earlier positive
differential occurs, so Proposition~\ref{fw:anss-pairing} makes
$d_{q+1}$ a derivation on this algebra. Let $\sigma$ denote the
algebraic specialization to $P$-cohomology. Differentiate the relation
and then apply $\sigma$. All terms containing $v_4$ specialize to zero,
and $d_{q+1}(v_4^p)=p v_4^{p-1}d_{q+1}(v_4)=0$. Hence
\begin{equation}
 \sigma\bigl(d_{q+1}(v_4)\bigr)h_{1,4}=0.
 \label{may:specialized-obstruction}
\end{equation}
Specialization is applied after the topological differential, to an
identity of cobar-cohomology classes; no spectrum realizing $\sigma$
is required.

It remains to prove that multiplication by $h_{1,4}$ is injective on
$H^{q+1,u_4+1}(P)$. Its degree shift is $(1,p^4)$, so its product
lies in degree $(q+2,u_5+1)$. An incoming May differential to that
product must have source degree $(q+1,u_5+1)$. The next lemma lists
the initial May monomials in the obstruction degree and in this
incoming-source degree.
We will use the digit-zero adjustment of Lemma~\ref{may:page} to
separate the surviving source filtrations from the product filtration.

In the alphabet of Equation~\eqref{may:low-alphabet}, the splitting
differential is
\begin{equation}
 da=xy,\qquad db=yz,\qquad dc=xb+az,
 \label{may:small-splitting}
\end{equation}
with the exterior order $x,a,c,y,b,z$ used for signs.

\begin{lemma}[The first target and all incoming-source intervals]
\label{may:first-list}
In bidegree $(s,w)=(q+1,u_4+1)$ the complete initial May page consists of
\begin{equation}
 acb B^{p-2},\qquad acybz B^{p-3}.
 \label{may:first-target-list}
\end{equation}
The first is not a splitting cycle and the second is a splitting cycle.
In bidegree $(q+1,u_5+1)$ the complete list is given in
Table~\ref{may:n5-table}.  The four entries of exterior length three or
five have no nonzero splitting cycle in their span.
\end{lemma}

\begin{proof}
We prove exhaustion simultaneously for $w=u_n+1$, with $n=4,5$.
Since $u_n+1<p^n$, the available exterior intervals are
$[a,b)$ with $0\leq a<b\leq n$, and the polynomial intervals cover
positions $1,\ldots,n-1$.  The target weight has base-$p$ digits
$(2,1,\ldots,1)$.  At position zero, $H_0\leq n<p$, so $H_0=2$ and
there is no carry into position one.  As $s=2p-1$ is odd, write
\[
 m=2D-1,\qquad L=p-D.
\]
There are exactly two exterior intervals starting at zero; at most
$n(n-1)/2$ others are available.  Hence $m\leq11$, $2\leq D\leq6$,
and
\begin{equation}
 p-6\leq L\leq p-2.
 \label{may:L-range}
\end{equation}
The lower bound on $D$ follows from $m$ being odd and at least two.

Let $c_r$ be the carry into position $r$, with $c_1=c_n=0$, and let
$B_r$ count polynomial intervals covering position $r$.  The largest
possible exterior coverage at a position is nine.  Inductively, the
quantity $H_r+B_r+c_r$ is at most $9+(p-2)+1=p+8<2p$;
therefore every carry is zero or one.  The digit equations are
\begin{equation}
 B_r=1-H_r-c_r+pc_{r+1},\qquad 1\leq r\leq n-1.
 \label{may:carry-equations}
\end{equation}
If there were no carry, then $B_r\leq1$ for every $r$, and
$L\leq\sum B_r\leq4$, contrary to $L\geq p-6=1003$.

Call a position high when $c_{r+1}=1$.  A high position has
$B_r\geq p-9$, and a nonhigh position has $B_r\leq1$.
There cannot be two blocks of high positions.  Indeed, choose high
positions $a<b$ in different blocks and a nonhigh position $c$ between
them.  An interval covering both $a$ and $b$ also covers $c$, whence
\[
 L\geq B_a+B_b-B_c\geq2(p-9)-1=2p-19>p-2.
\]
Here $2p-19=1999$ and $p-2=1007$.  This contradicts
\eqref{may:L-range}.  Thus the high positions form a single block
$[l,r]$ with $1\leq l\leq r\leq n-2$.
At the next position \eqref{may:carry-equations} reads
$B_{r+1}=-H_{r+1}$, so
\begin{equation}
 B_{r+1}=H_{r+1}=0.
 \label{may:gap-after-block}
\end{equation}
At the beginning of the block, $B_l=p+1-H_l\leq L$, giving
$H_l\geq D+1$.  If $l<r$, its end gives
$B_r=p-H_r\leq L$, hence $H_r\geq D$.

A singleton block is impossible.  Every exterior interval covering its
position $l$ must end at $l+1$ by \eqref{may:gap-after-block}.  When
$l=1$ there are only two such intervals.  When $l>1$, at most one begins
before $l$, since $H_{l-1}\leq1$, and at most one begins at $l$.
Thus in either case $H_l\leq2$, contradicting $H_l\geq D+1\geq3$.
We henceforth have $l<r$.

If $l>1$, at most two exterior intervals cover both $l$ and $r$: they
must end at $r+1$, and at most one starts before $l$, while at most one
starts at $l$.  At least one of the two intervals starting at zero
misses both positions, since $H_{l-1}\leq1$.  Consequently
\[
 H_l+H_r
 =m+\#\{\text{intervals meeting both}\}
       -\#\{\text{intervals meeting neither}\}
 \leq m+1.
\]
This contradicts $H_l+H_r\geq2D+1=m+2$.  Therefore $l=1$.
There are now at most two exterior intervals covering both endpoints,
namely $[0,r+1)$ and $[1,r+1)$.  Equality in the same counting identity
is forced: both are present, every exterior interval meets position $1$
or position $r$, and
\[
 H_1=D+1,\qquad H_r=D.
\]
It follows that $B_1=B_r=L$.  Every polynomial interval therefore covers
both positions, and by \eqref{may:gap-after-block} every one is exactly
$[1,r+1)$.  No exterior or polynomial interval can occur after the block.
The carry supplies the target digit at $r+1$; all further digits would
be zero.  Since the required remaining digits are all one, necessarily
$r=n-2$.  Thus
\begin{equation}
 \text{polynomial factor }b_{n-2,0}^{p-D},\qquad
 H_0=2,\quad H_1=D+1,\quad
 H_2=\cdots=H_{n-2}=D,
 \quad H_{n-1}=0.
 \label{may:forced-profile}
\end{equation}

It remains to list a small, now forced set of exterior intervals.
For $n=4$, the intervals $[0,3)$ and $[1,3)$ are mandatory.  The other
start-zero interval must be $[0,2)$; $[0,1)$ would meet neither endpoint.
The only optional intervals are $[1,2)$ and $[2,3)$.  The coverage
requirements in \eqref{may:forced-profile} require both or neither.
This is exactly \eqref{may:first-target-list}.

For $n=5$, the mandatory intervals are $[0,4)$ and $[1,4)$.
The remaining start-zero interval is $[0,t)$ with $t=2$ or $3$.
The only optional intervals are $[1,2),[1,3),[2,4),[3,4)$;
$[2,3)$ meets neither endpoint and is forbidden.  Write their inclusion
indicators as $e_{12},e_{13},e_{24},e_{34}\in\{0,1\}$, and let their
sum be $o$.  The three required coverage equations are
\begin{equation}
 e_{12}+e_{13}=o/2,\qquad
 e_{24}+e_{34}=o/2,\qquad
 e_{13}+e_{24}+[t=3]=o/2,
 \label{may:four-indicators}
\end{equation}
where $o\in\{0,2,4\}$.  For $t=2$ the solutions are no optional
intervals, the pair $\{12,24\}$, the pair $\{13,34\}$, and all four.
For $t=3$ the only solution is the pair $\{12,34\}$.  This proves the
exhaustive Table~\ref{may:n5-table}, including its length-seven entry.

\begin{table}[htbp]
\centering
\begin{tabular}{c l l}
\toprule
$m$ & Exterior intervals (exclusive right endpoints) & Polynomial factor\\
\midrule
$3$ & $02,04,14$ & $b_{3,0}^{p-2}$\\
$5$ & $03,04,12,14,34$ & $b_{3,0}^{p-3}$\\
$5$ & $02,04,13,14,34$ & $b_{3,0}^{p-3}$\\
$5$ & $02,04,12,14,24$ & $b_{3,0}^{p-3}$\\
$7$ & $02,04,12,13,14,24,34$ & $b_{3,0}^{p-4}$\\
\bottomrule
\end{tabular}
\caption{Every initial May monomial in bidegree $(q+1,u_5+1)$.
The symbol $ab$ in this table denotes the interval $[a,b)$, not a
product of the letters in \eqref{may:low-alphabet}.}
\label{may:n5-table}
\end{table}

For completeness, the splitting calculation is also explicit.
The length-three entry has, for example, a nonzero component on
$01,04,12,14$.  For the three length-five entries in their displayed
order, take the following four rows, ordering exterior intervals
lexicographically:
\[
\begin{array}{c|rrr}
01,04,12,13,14,34&1&-1&0\\
02,03,12,14,24,34&1&0&1\\
02,04,12,13,24,34&0&1&1\\
02,04,12,14,23,34&-1&-1&1
\end{array}
\]
The common polynomial factor is $b_{3,0}^{p-3}$.  The minor using the
first, second, and fourth rows is $3$, which is nonzero in
$\mathbb F_{1009}$.  Thus these three images are linearly independent.
Finally, \eqref{may:small-splitting} shows that $acb$ has a nonzero
splitting differential and that $acybz$ is closed: every term in the
latter differential repeats an exterior factor.  This finishes the
lemma.
\end{proof}

\begin{proposition}[An injective multiplication]
\label{may:injectivity}
Multiplication by the primitive class $h_{1,4}$ is injective on actual
cobar cohomology in bidegree $(q+1,u_4+1)$:
\[
 H^{q+1,u_4+1}(P)\xrightarrow{\ h_{1,4}\ }
 H^{q+2,u_5+1}(P).
\]
\end{proposition}

\begin{proof}
Use the filtration in Lemma~\ref{may:page} with the explicit choice
$M=3p$.  If the source cohomology is nonzero, any nonzero class has
leading detector a nonzero scalar multiple of
$acybz B^{p-3}$, by Lemma~\ref{may:first-list}.  Its filtration is
\[
 F_0=3p^2-9p+13+M(2+p(p-3)).
\]
The cocycle $[t_1^{p^4}]$ represents $h_{1,4}$, with degree, weight, and
filtration shifts $(1,p^4,1)$.  The proposed leading product therefore
has filtration
\begin{equation}
 F_*=3p^2-9p+14+M(2+p(p-3)).
 \label{may:product-filtration}
\end{equation}

Every possible incoming May differential to this product has source
cohomological degree $q+1$ and weight $u_5+1$.
Table~\ref{may:n5-table} is its complete initial-page source list.
The length-three and length-five source filtrations are respectively
\[
 5p^2-10p+15+M(2+p(p-2)),\qquad
 5p^2-15p+19+M(2+p(p-3));
\]
all three length-five sources have the latter filtration.
The length-three and length-five entries disappear under $d_1$, since
their displayed splitting maps are injective.  None can be a splitting
boundary for the product itself, because its polynomial factor is a
power of $B=b_{2,0}$, whereas their polynomial factor is a power of
$b_{3,0}$.  The length-seven source has filtration
\[
 F_7=5p^2-20p+23+M(2+p(p-4)).
\]
For $M=3p$ the difference is
\[
 F_7-F_*=2p^2-11p+9-pM=-p^2-11p+9<0.
\]
As all May differentials lower filtration, this last source cannot hit
the proposed product on any later page.  There are no other incoming
sources.

This argument does not require $B$ to represent a permanent class by
itself.  Choose an actual cocycle for the given nonzero source
cohomology class and multiply that entire cocycle by
$[t_1^{p^4}]$.  The result is an actual cocycle.  Its leading
associated-graded cup product is the nonzero monomial
$acybz B^{p-3}h_{1,4}$, since the new exterior letter is distinct from
all the old ones.  An actual cocycle has no outgoing May differential;
its leading detector could disappear only by becoming a boundary,
which the preceding incoming-source calculation excludes at every
page.  Its product cohomology class is therefore nonzero.
\end{proof}

\begin{proposition}[The first Adams--Novikov differential]
\label{may:first-differential}
In the Adams--Novikov spectral sequence for $Y$, one has
$d_{q+1}(v_4)=0$.
\end{proposition}

\begin{proof}
Equation~\eqref{may:specialized-obstruction} and
Proposition~\ref{may:injectivity} give
$\sigma(d_{q+1}(v_4))=0$. Lemma~\ref{may:target-bound} identifies
this specialization with the original target group. Hence
$d_{q+1}(v_4)=0$.
\end{proof}

\subsection{The remaining targets}

\begin{proposition}[All later obstruction candidates]
\label{may:remaining-list}
For $k\geq2$ the complete initial May page in
\eqref{may:obstruction-degrees} is Table~\ref{may:remaining-table}; all
other values of $k$ give an empty page.  Put
$N_0=(p-1)^2$ and $N=p(p-1)-2$.
\begin{table}[htbp]
\centering
\begin{tabular}{c l}
\toprule
$k$ & Initial May monomials\\
\midrule
$2$ & $xac B A^{p-2}C^{p-2}$,
       $\quad xacyz A^{p-2}C^{p-2}$\\[2pt]
$p-1$ & $y B^3 A^{N_0-3}$,
         $\quad b B^2 A^{N_0-2}$,
         $\quad ybz B A^{N_0-2}$\\[2pt]
$p$ & $cyb A^{N+1}$,
       $\quad ayb B A^N$\\
\bottomrule
\end{tabular}
\caption{Every initial May monomial for the later $v_4$ obstruction
bidegrees.}
\label{may:remaining-table}
\end{table}
\end{proposition}

\begin{proof}
Lemma~\ref{may:target-bound} allows only positions $0,1,2$ and
$2\leq k\leq p+2$.  Thus \eqref{may:three-position-weight} applies.
There are six exterior letters; $m$ is odd, so write
\[
 m=2f+1,\quad 0\leq f\leq2,\qquad L=(p-1)k-f.
\]
Reducing the weight modulo $p$ gives
$H_0\equiv1+k\pmod p$, with $0\leq H_0\leq3$.
The only possibilities are $k=2$ or $k=p+d$ with
$-1\leq d\leq2$.

When $k=2$, $H_0=3$ forces all of $x,a,c$.  The other exterior letters
are an even subset of $\{y,b,z\}$.  Put
$T=H_2+n_B+n_C$.  Substitution in
\eqref{may:three-position-weight} gives
\begin{equation}
 pT-n_C+H_1=p^2-p+3+f.
 \label{may:k2-equation}
\end{equation}
For the four possible exterior subsets the value
$e=H_1-3-f$ is $-2$ or $-1$.  Therefore
$n_C\equiv e\pmod p$.  Since
$0\leq n_C\leq L\leq2p-3$, the only possible lift is
$n_C=p+e$: the lift $e$ is negative, and $2p+e\geq2p-2>L$.
Equation~\eqref{may:k2-equation} now gives $T=p$.
All the remaining counts follow as displayed:
\[
\begin{array}{c|c|c|c|c|c|c}
\text{extra exterior letters}&f&H_1&H_2&n_C&n_B&n_A\\ \hline
\varnothing&1&2&1&p-2&1&p-2\\
y,b&2&4&2&p-1&-1&p-2\\
y,z&2&3&2&p-2&0&p-2\\
b,z&2&3&3&p-2&-1&p-1
\end{array}
\]
The negative counts exclude the second and fourth rows.  The first and
third are exactly the $k=2$ row of Table~\ref{may:remaining-table}.

Now let $k=p+d$, $-1\leq d\leq2$.  Then $H_0=1+d$, and the weight
equation becomes
\begin{equation}
 pT-n_C+H_1=(2-d)p+2+d+f.
 \label{may:pd-equation}
\end{equation}
Since $n_C\leq T$ and $H_1\geq0$, it follows that
\[
 (p-1)T\leq(2-d)(p-1)+(4+f).
\]
Here $4+f\leq6<p-1$, so $T\leq2-d\leq3$ and $n_C\leq3$.
Rearranging \eqref{may:pd-equation} now gives
\[
 p(T-(2-d))=2+d+f+n_C-H_1.
\]
The right side lies between $-3$ and $9$, strictly between $-p$ and
$p$.  Both sides therefore vanish, proving the exact equations
\begin{equation}
 H_2+n_B+n_C=2-d,\qquad H_1=2+d+f+n_C.
 \label{may:pd-exact}
\end{equation}

For $d\geq1$, these require
$3+f\leq H_1\leq\min(2f+1,4)$, impossible for each of
$f=0,1,2$.  For $d=0$, exactly one of $x,a,c$ occurs, so $H_1\leq3$.
Equation~\eqref{may:pd-exact} forces $f=1$, $n_C=0$, and $H_1=3$.
The only exterior triples are $ayb$ and $cyb$.  Their $H_2$ values are
respectively $1$ and $2$, giving the two $k=p$ monomials in the table.

For $d=-1$, only $y,b,z$ can occur.  If $m=1$, the choices $y,b,z$
have $(H_1,H_2)$ equal to $(1,0),(1,1),(0,1)$.
Equation~\eqref{may:pd-exact} excludes $z$ and gives respectively
$(n_C,n_B)=(0,3),(0,2)$ for the other two.  If $m=3$, the exterior
factor is $ybz$, with $f=1$, $H_1=H_2=2$; the same equations give
$(n_C,n_B)=(0,1)$.  Since $L=N_0-f$, these are exactly the three
$k=p-1$ monomials.  Every nonnegative count and every exterior subset
has now been considered.
\end{proof}

The first two rows of Table~\ref{may:remaining-table} use only the first
three or two coordinates, respectively. Every monomial in its last row
has a large common power of $A$. The following comparisons turn these
two observations into surjections on actual cobar cohomology, which are
the inputs to the remaining obstruction arguments.

\begin{lemma}[Two comparison surjections]
\label{may:comparison}
The following initial-page tests give surjections on actual cobar
cohomology.
\begin{enumerate}
 \item If every monomial of \eqref{may:initial-page} in bidegree
 $(s,w)$ uses only generators with $i\leq r$, then
 $H^{s,w}(P_r)\longrightarrow H^{s,w}(P)$ is surjective.
 \item Put $A=b_{1,0}$. If every initial-page monomial in bidegree
 $(s,w)$ contains $A^a$, then multiplication by the actual cocycle
 representing $A^a$ gives a surjection
 \[
 H^{s-2a,w-pa}(P)\longrightarrow H^{s,w}(P).
 \]
\end{enumerate}
These conclusions also transport through additive comparisons that
identify the source and target cohomology groups and commute with the
indicated maps.
\end{lemma}

\begin{proof}
The subalgebra map includes a subset of the exterior and polynomial
generators in \eqref{may:initial-page}; its initial-page map is therefore
injective in every bidegree. Likewise, multiplication by the polynomial
$A^a$ is injective in every bidegree on that page. It is induced by a
filtered cochain map: $t_1$ is primitive, and
\begin{equation}
 \sum_{i=1}^{p-1}\frac{\binom pi}{p}
       [t_1^i\mid t_1^{p-i}]
 \label{may:A-cocycle}
\end{equation}
is an actual cocycle, obtained by taking the integral divided reduced
coproduct and then reducing modulo $p$. It has leading class $A$.
For the second comparison shift the source cohomological degree,
weight, and auxiliary filtration by $2a$, $pa$, and $ap(1+M)$.

Take the filtered mapping cone of either cochain map. The long exact
sequence of associated-graded cohomology, together with injectivity in
\emph{all} bidegrees, identifies its initial cohomology page with the
cokernel of the specified page map; there is no extra kernel from the
next degree. That cokernel is zero in the stated target bidegree.
Finite filtration therefore gives zero cone cohomology there. The
ordinary cone long exact sequence gives the required surjection.
Transport through the stated commuting comparisons preserves that
surjectivity.
\end{proof}

Section~\ref{frob:section} removes the first two rows of
Table~\ref{may:remaining-table} in actual cobar cohomology.  The last
row is treated in Section~\ref{top:section} by a homotopy argument with an explicit
Adams--Novikov page bound.
